\scanpage{55}
\mspage{- 25 -}

n\textsuperscript{o} 276

\uline{couple} $(\hat{\mathrm A},\mathrm m)$ \uline{est un anneau de Zariski}.

(b) - \uline{Pour tout} $\mathrm A$\uline{-module} $\mathrm E$ \uline{de type fini}, \uline{l'application canonique}
\[
\hat{\mathrm A}\underset{\mathrm A}{\handwritten{\otimes}}\mathrm E\handwritten{\longrightarrow}\hat{\mathrm E}
\]
\uline{est bijective}.

(c) - \uline{Le couple} $(\mathrm A,\hat{\mathrm A})$ \uline{est plat}.

Soit $\mathrm u:\mathrm E\handwritten{\longrightarrow}\mathrm F$ un homomorphisme de $\mathrm A$-modules de type fini. Comme on a $\mathrm u(\mathrm m^{\mathrm p}\mathrm E)=\mathrm m^{\mathrm p}\mathrm u(\mathrm E)$, le groupe topologique $\mathrm u(\mathrm E)$ est quotient du groupe topologique $\mathrm E$ ; or la topologie $\mathrm m$-adique de $\mathrm u(\mathrm E)$ est induite par celle de $\mathrm F$ (Théorème 1) ; par conséquent, $\mathrm u$ est un \uline{homomorphisme topologique} de $\mathrm E$ dans $\mathrm F$. Il résulte de là et du lemme 1 du \handwritten{\S}1, n\textsuperscript{o}7, que le foncteur $\mathrm E\handwritten{\to}\hat{\mathrm E}$ transforme toute suite exacte de $\mathrm A$-modules \uline{de type fini} en une suite exacte.

Pour établir l'assertion (b) observons qu'on a une suite exacte
\[
\mathrm L_1\handwritten{\longrightarrow}\mathrm L_{\unsure{\mathrm o|0}}\handwritten{\longrightarrow}\mathrm E\handwritten{\longrightarrow}0
\]
avec $\mathrm L_{\mathrm o}$ et $\mathrm L_1$ libres et de type fini puisque $\mathrm A$ est noethérien. D'après ce qui précède et les propriétés élémentaires des produits tensoriels, il vient donc un diagramme commutatif de suites exactes

\diagram{\begin{tikzcd}[ampersand replacement=\&]
\hat{\mathrm A}\handwritten{\otimes}\mathrm L_1 \handwritten{\arrow[r]}\handwritten{\arrow[d]} \& \hat{\mathrm A}\handwritten{\otimes}\mathrm L_{\mathrm o} \handwritten{\arrow[r]}\handwritten{\arrow[d]} \& \hat{\mathrm A}\handwritten{\otimes}\mathrm E \handwritten{\arrow[r]}\handwritten{\arrow[d]} \& 0 \\
\hat{\mathrm L}_1 \handwritten{\arrow[r]} \& \hat{\mathrm L}_{\mathrm o} \handwritten{\arrow[r]} \& \hat{\mathrm E} \handwritten{\arrow[r]} \& 0
\end{tikzcd}}
% The nine arrows of the diagram are drawn in by hand.

De sorte que pour établir que le troisième homomorphisme vertical est bijectif il suffit d'établir que les deux premiers le sont : mais si $\mathrm E$ est isomorphe à $\mathrm A^{\mathrm n}$ il est clair que $\hat{\mathrm E}$ est isomorphe à $\hat{\mathrm A}^{\mathrm n}\unsure{.}=\hat{\mathrm A}\underset{\mathrm A}{\handwritten{\otimes}}\mathrm A^{\mathrm n}$, d'où le résultat.

Il suit de là que le foncteur $\mathrm E\handwritten{\to}\hat{\mathrm A}\underset{\mathrm A}{\handwritten{\otimes}}\mathrm E$ est exact sur la catégorie des $\mathrm A$-modules de type fini ; le $\mathrm A$-module $\hat{\mathrm A}$ est donc plat. Comme de plus l'homomorphisme $\mathrm E\handwritten{\to}\hat{\mathrm A}\underset{\mathrm A}{\handwritten{\otimes}}\mathrm E=\hat{\mathrm E}$ est injectif si $\mathrm E$ est de type fini - ceci parce que $\mathrm E$ est séparé pour la topologie \quad -adique en vertu du théorème 2-

% archive sheet number: 54
